\subsection{Restriction of scalars}

PROPOSITION 7. Let A, B be two rings and \(\rho\) a homomorphism of A to B. Let E be a faithfully flat right B-module. For \(\rho^{*}(\mathrm{E})\) to be a flat (resp. faithfully flat) right A-module, it is necessary and sufficient that B be a flat (resp. faithfully flat) right A-module.

Applying Proposition 4 of no. 2 with B, A, E, B in place of R, S, E, F respectively, and the right A-module structure on B being defined by \(\rho\) , it is seen that

B is a flat (resp. faithfully flat) A-module if and only if \(E \otimes_{B} B = \rho^{*}(E)\) is a flat (resp. faithfully flat) A-module.

\textbf{Remarks}\par

\begin{enumerate}
\def\labelenumi{(\arabic{enumi})}
\item
  Proposition 7 shows that, for B to be a faithfully flat A-module, it is sufficient that there exist one faithfully flat B-module which is also a faithfully flat A-module.
\item
  Let A, B, C be three rings and \(\rho: A \rightarrow B\) , \(\sigma: B \rightarrow C\) two ring homomorphisms. Proposition 7 shows that if C is a faithfully flat B-module and B a faithfully flat A-module, then C is a faithfully flat A-module. If C is a faithfully flat B-module and a faithfully flat A-module, then B is a faithfully flat A-module (taking the modules as right modules, to fix the ideas). On the other hand B and C may be faithfully flat A-modules without C being a faithfully flat B-module (Exercise 7).
\end{enumerate}

